% §8 Concluding remarks --- paper-only, written 2026-09-12 for the scope §§1-7 (v0.1 stops at §7,
% the author's decision of that day). On Paper I's model: the summary, the scope, one
% methodological remark on what the formalization changed in the printed proofs (each item is a
% blueprint CHANGED marker in chapters 5-7), and the open questions --- stated as open, not as
% promised (hub RELEASES.md rule 1). No blueprint node is shared here.

\section{Concluding remarks}\label{sec:conclusions}

This paper weakened one axiom of \sscite{pauwels1995extended}, with two smaller changes
recorded in Remark~\ref{rem:provenance-pauwels}. Recursivity, the one-parameter semigroup in
scale, carries the tacit assumption that an increment depends on the scale difference
alone, the stationarity of the increments in the scale parameter; it was replaced by the
two-parameter hemigroup. Under the remaining axioms
the admissible measurement families on the line are exactly the symmetric self-decomposable
laws, each determined by a Gaussian coefficient and a nonincreasing displacement profile,
with the gauge and the exponent unique up to one normalization
(Theorem~\ref{thm:main-characterization}). The semigroup case is the symmetric stable
family with $\alpha \le 2$ (Corollary~\ref{cor:semigroup-case}), so the heavy tails of the
stable members below $\alpha = 2$ are a consequence of stationarity and not of
covariance. Positivity is the narrowing axiom: without it the hemigroup axioms admit an
infinite-dimensional affine family through the Gaussian
(Proposition~\ref{prop:no-positivity-no-classification}, whose construction is the one
step of this paper not machine-checked), and with it the class is a convex cone in which the
Gaussian scale space is the extreme ray without jumps. The kernels from the origin are
absolutely continuous and unimodal (Proposition~\ref{prop:kernel-regularity}). No truncation of the
Gaussian is admissible, since its transform changes sign (Lemma~\ref{lem:nonvanishing}).

On scope: this paper delivers the classification on the line, in the $L^1$ signal model, and
nothing more. It makes no empirical claims about natural signals. The geometry of the cone,
the families at its corners, the scale evolution in generator form, locality, non-creation
of local extrema, the exact implementation of the cascade, higher dimension and the
discrete covariance strata are outside it.

In relation to the axiom systems of \S\ref{sec:introduction}: the Gaussian axiomatics,
from Iijima's conditions \sscite{iijima1962basic,weickert1999linear} to the non-creation
and non-enhancement principles \sscite{koenderink1984structure,lindeberg1994scale}, keep
the semigroup and add a requirement on how structure may change across scale. Within the
class characterized here, the semigroup selects the stable slice and such a requirement
selects its end $\alpha = 2$. The extended class of \sscite{pauwels1995extended} and the
$\alpha$-scale spaces of \sscite{duits2004axioms} are that slice, and the Poisson scale
space of \sscite{felsberg2001scale} is its member $\alpha = 1$. The Bessel scale space of
\sscite{burgeth2005bessel}, composed along the range, is the variance-gamma member, with jumps of
every size and every moment finite. The time-causal classification of
\sscite{fagerstrom2026hemigroup} is the corresponding theorem under the Laplace transform,
with a drift term where the Gaussian term is here; the axiom that differs, causality there
and reflection symmetry here, is what excludes the Gaussian on the half-line and admits it
on the line.

One methodological remark. Several of the proofs printed above are not the ones first written, and in
each case the machine-checked route is shorter or assumes less. The null-array limit of
Theorem~\ref{thm:increments-levy} extracts the representation of an increment through a
truncation inequality proved in a few lines, and never uses the closure of the class under
pointwise limits. The analysis direction of Lemma~\ref{lem:selfdecomposable-exponents} first
argued through the convexity of a tail function and a right derivative. The checked proof
defines the profile as a supremum of dyadic difference quotients, which is nonincreasing by
construction, and uses no convexity at all. The continuity axiom in the sufficiency half of
Theorem~\ref{thm:main-characterization} was first verified by an $\varepsilon/3$ argument on
a compact set carrying most of the mass. The checked proof bounds the distance of a
convolution from its argument by the average of the modulus of continuity of translation,
and needs no compact set and no density argument. The strict positivity of the admissible
exponents (Proposition~\ref{prop:strict-positivity}) first covered the half-line by
intervals and passed from almost everywhere to everywhere. The checked proof reads a
dilation-invariant tail integral at one point. And the bound $\alpha \le 2$ of the
semigroup case is the three-line quadratic growth estimate rather than a cited one. The net
effect is that the characterization rests on the two cited facts of
\S\ref{sec:machine-checked}, and that the printed proofs are the verified ones. The two
further clauses the classical route uses, the closure of the class under pointwise limits
and the existence half of the \Levy--Khintchine representation, are used nowhere.

Four questions are left open. First, the construction in
Proposition~\ref{prop:no-positivity-no-classification} (the signed family built from an
even integrable transform) is the one step of this paper that is not machine-checked. A
route through the convolution operators has been estimated and not carried out. Second, the
non-positive hemigroup class is shown to admit no description in the form of the theorem,
and it is not described: what it contains beyond the affine families through the admissible
exponents is not known. Third,
the closed form of the Mat\'ern member's kernel is cited in \S\ref{sec:axioms} and not
claimed, and a self-contained derivation of the densities of the members with a completely
monotone profile would make it a statement of the paper.
Fourth, the classification is on the line. Which of it survives rotation invariance in
dimension $d$, and which of it survives covariance under a discrete group of dilations only,
are the questions the hemigroup form of the axioms raises next. The probability classes
that would answer them are known: the self-decomposable laws and self-similar additive
processes in $\RR^d$ (\sscite{sato1991selfsimilar}), and for a discrete group the
semi-selfdecomposable laws and semi-selfsimilar processes of Maejima and Sato
(\sscite{maejima1999semi}, Def.~1, p.~348, and Thm.~6, p.~360). What is open is the
derivation from measurement axioms, and its formalization.
